\subsection{Structural function of \texorpdfstring{\(K\) and \(\alpha\)}{K and alpha} in exceptional incidence}

The link between numerical determination and exceptional realization
admits an operational description. In the arithmetic reading, \(K\)
acts in the normalization of regional records. In the incidence
reading, its components select the support
\[
 B_K=\{j:K_j\geq729\}=\{4,6,9,10\}.
\]
The same ordered record therefore participates in two specified
functions: a positional composition and an incidence selection.
Its rational realization preserves the record when base,
length, and origin are fixed, while passage to the support retains the information
needed for the combinatorial flag.

The signed record preceding normalization of \(\alpha\) supplies
the hexad \(H_\alpha^-\). The intersections of this hexad with the
regional supports and APP incidence determine the marked
point. That point individualizes the norm-\(54\) vector used
in the lattice neighbor. Thus the orientation entering arithmetic
normalization acts again in geometric selection.
The relation does not require deriving a lattice from a decimal
approximation: it uses the signed data preceding its evaluation.

The exceptional realization successively preserves different types of
relation. The code preserves linear combinations; the design preserves
incidences between supports; gluing preserves discriminant classes; the
lattice preserves its bilinear form; the vertex algebra preserves grading
and the corresponding products. Their automorphism groups
are defined with respect to these structures. The appearance of the Monster at
the end of the construction expresses a property of the algebra
\(V^\natural\) obtained through the cited procedure, and Moonshine
relates the graded traces of that action to modular functions.

This hierarchy gives a concrete mathematical meaning to the
linking function attributed to \(K\) and \(\alpha\). The former is a
reversible integral record with positional and incidence
realizations; the latter is the arithmetic-compatibility coordinate
whose oriented record also enters the marked incidence.
The definitions acquire content through these operations and the
identities presented in this chapter. The incidence development
is a substantive part of the article's argument: it shows which relations
persist when the construction is no longer evaluated as a real
coordinate and is instead realized as an algebraic structure.
